%%% ---
%%% Packages
%%% --- 

% General
\usepackage{amsmath,amsthm,amssymb,amsfonts, mathabx, tikz-cd}
\usepackage{commath}
\usepackage{extarrows}

% Graphics
\usepackage{tikz}
\usetikzlibrary{decorations.pathreplacing}
\usepackage[dvipsnames]{xcolor}
\usepackage{tabularx}
\usepackage{longtable}
\usepackage{array}
\usepackage{faktor}

% Polygraphic pictures
\usepackage{catex}
\usepackage[all]{xy}

% Lists and floats
\usepackage{enumitem}
\usepackage{float}

% Structure / appendix
\usepackage{titletoc}
\usepackage[titletoc]{appendix}

% Layout
\usepackage{geometry}

% Hyperlinks
\usepackage[hyperfootnotes=false,draft=false,pdftex, colorlinks=true, linkcolor=black, citecolor=black, urlcolor=black, hypertexnames=false]{hyperref}

%% Margins
\geometry{marginparwidth=21mm, marginparsep=1mm, left=25mm, right=25mm, top=25mm, bottom=25mm}
\usepackage{adjustbox}

%Romanian sch
\usepackage{combelow}
%%% --- 
%%% General notation
%%% --- 

\numberwithin{equation}{section}

\newcommand{\defeq}{\,:=\,}
\newcommand{\eq}{\,=\,}
\newcommand{\less}{\,<\,}
\newcommand{\greater}{\,>\,}

\newcommand\eqdef{\mathrel{\overset{\makebox[0pt]{\mbox{\normalfont\tiny\sffamily def}}}{\eq}}}

\newcommand\eqabove[1]{\mathrel{\overset{\makebox[0pt]{\mbox{\normalfont\footnotesize #1}}}{=}}}
\newcommand\leqabove[1]{\mathrel{\overset{\makebox[0pt]{\mbox{\normalfont\footnotesize #1}}}{\leq}}}

\newcommand{\ms}[1]{\mathsf{#1}}
\newcommand{\vect}[1]{\vec{#1}\,}
\renewcommand{\norm}[1]{\lVert #1 \rVert}


%%% --- 
%%% Formatting and visual tweaks
%%% --- 

\newcommand{\textbfit}[1]{\textbf{\textit{#1}}}
\newcommand\Simeq{\;\simeq\;}

\newcommand{\overbar}[1]{\mkern 1mu\overline{\mkern-1mu#1\mkern-1mu}\mkern 1mu}
\newcommand{\highoverline}[1]{\overbar{\vphantom{#1^{\rule{0pt}{0.3ex}}}#1}}


%%% --- 
%%% Categories and functors
%%% --- 

\newcommand{\LDl}{\ms{uLD}_{\mathrm{lax}}}
\newcommand{\LD}{\ms{uLD}}
\newcommand{\Monl}{\ms{uMon}_{\mathrm{lax}}}
\newcommand{\uFrob}{\ms{uFrob}}

\newcommand{\Free}{\mathsf{Free}_{\fS}}
\newcommand{\FreeC}{\mathsf{Free}_{{\operatorname{Ob(\cC)}}}}  

\newcommand{\rotamp}[2]{\mathbin{\rotatebox[origin=c]{180}{$#1\&$}}}
\newcommand{\parll}{{\mathpalette\rotamp\relax}}
\DeclareMathOperator{\partimes}{{\parll}}
\newcommand{\ot}{\otimes}

\newcommand{\tim}{\times}

\newcommand{\Hot}{H^{\ot}}
\newcommand{\Lot}{L^{\ot}}
\newcommand{\Rot}{R^{\,\ot}}

\newcommand{\Hpar}{H^{\partimes}}
\newcommand{\Lpar}{L^{\partimes}}
\newcommand{\Rpar}{R^{\,\partimes}}

\newcommand{\id}{{\operatorname{id}}}
\newcommand{\idC}{{\operatorname{id}_{\cC}}}

\newcommand{\Hom}[2]{\operatorname{Hom}_{\cC}(#1,#2)}

\newcommand{\dom}{{\operatorname{dom}}}
\newcommand{\codom}{{\operatorname{codom}}}
\newcommand{\Ob}{{\operatorname{Ob}}}

%%% --- 
%%% Coherence morphisms
%%% --- 

\newcommand{\ao}{\alpha}
\newcommand{\aoi}{\alpha^{-1}}
\newcommand{\lo}{\lambda}
\newcommand{\loi}{\lambda^{-1}}
\newcommand{\ro}{\rho}
\newcommand{\roi}{\rho^{-1}}

\newcommand{\aoc}{a}
\newcommand{\aoic}{a^{-1}}
\newcommand{\loc}{l}
\newcommand{\loic}{l^{-1}}
\newcommand{\roc}{r}
\newcommand{\roic}{r^{-1}}

\newcommand{\ap}{\highoverline{\alpha}}
\newcommand{\api}{\highoverline{\alpha}^{-1}}
\newcommand{\lp}{{\highoverline{\lo}}}
\newcommand{\lpi}{{\highoverline{\loi}}}
\newcommand{\rp}{{\overline{\ro}}}
\newcommand{\rpi}{{\highoverline{\roi}}}

\newcommand{\apc}{\overline{a}}
\newcommand{\apic}{\highoverline{a}^{-1}}
\newcommand{\lpc}{{\overline{l}}}
\newcommand{\lpic}{{\highoverline{\loic}}}
\newcommand{\rpc}{{\overline{\roc}}}
\newcommand{\rpic}{{\highoverline{\roic}}}

\newcommand{\distl}{\delta^l}
\newcommand{\distr}{\delta^r}

\newcommand{\distlc}{d^l}
\newcommand{\distrc}{d^r}

%%% --- 
%%% Arrows
%%% --- 

\newcommand{\slongrightarrow}[1]{\,\longrightarrow{#1}\,}
\newcommand{\sxlongrightarrow}[1]{\,\xlongrightarrow{#1}\,}


%%% --- 
%%% Frequently used symbols
%%% --- 

\newcommand{\oLD}{\highoverline{\ms{LD}}}
\newcommand{\Ord}{\mathbf{Ord}} 
\newcommand{\fN}{\mathfrak{N}}
\newcommand{\fS}{S}
\newcommand{\fY}{\mathfrak{Y}}
\newcommand{\fZ}{\mathfrak{Z}}

\newcommand{\Neq}{\mathrel{\scalebox{1.1}{$\neq$}}}
\newcommand{\NNez}{\NN_{0}}

\newcommand{\e}{\emptyset}


%%% --- 
%%% Standard symbols
%%% --- 

\def\cA{\mathcal A}
\def\cB{\mathcal B}
\def\cC{\mathcal C}
\def\cD{\mathcal D}
\def\cE{\mathcal E}
\def\cF{\mathcal F}
\def\cG{\mathcal G}
\def\cH{\mathcal H}
\def\cI{\mathcal I}
\def\cJ{\mathcal J}
\def\cK{\mathcal K}
\def\cL{\mathcal L}
\def\cM{\mathcal M}
\def\cN{\mathcal N}
\def\cO{\mathcal O}
\def\cP{\mathcal P}
\def\cQ{\mathcal Q}
\def\cR{\mathcal R}
\def\cS{\mathcal S}
\def\cT{\mathcal T}
\def\cU{\mathcal U}
\def\cV{\mathcal V}
\def\cW{\mathcal W}
\def\cX{\mathcal X}
\def\cY{\mathcal Y}
\def\cZ{\mathcal Z}

\def\AA{\mathbb A}
\def\BB{\mathbb B}
\def\CC{\mathbb C}
\def\DD{\mathbb D}
\def\EE{\mathbb E}
\def\FF{\mathbb F}
\def\GG{\mathbb G}
\def\HH{\mathbb H}
\def\II{\mathbb I}
\def\JJ{\mathbb J}
\def\KK{\mathbb K}
\def\LL{\mathbb L}
\def\MM{\mathbb M}
\def\NN{\mathbb N}
\def\OO{\mathbb O}
\def\PP{\mathbb P}
\def\QQ{\mathbb Q}
\def\RR{\mathbb R}
\def\SS{\mathbb S}
\def\TT{\mathbb T}
\def\UU{\mathbb U}
\def\VV{\mathbb V}
\def\WW{\mathbb W}
\def\XX{\mathbb X}
\def\YY{\mathbb Y}
\def\ZZ{\mathbb Z}
\def\Si{\Sigma}
\def\wSi{\widehat{\Sigma}}


%%% --- 
%%% Theorem environments
%%% --- 

\theoremstyle{plain}
\newtheorem{theorem}{Theorem}[section]
\newtheorem*{maintheorem*}{Main Theorem}
\newtheorem{lemma}[theorem]{Lemma}
\newtheorem{proposition}[theorem]{Proposition}              
\newtheorem{corollary}[theorem]{Corollary}        

\newtheorem*{cohthm}{Linearly Distributive Coherence Theorem}
\newtheorem*{lapthm}{Lax Monoidal Coherence Theorem}

\theoremstyle{definition} 
\newtheorem{definition}[theorem]{Definition}

\theoremstyle{remark}
\newtheorem{remark}[theorem]{Remark}
\newtheorem{example}[theorem]{Example}
\newtheorem{construction}[theorem]{Construction}
\newtheorem{setup}[theorem]{Setup}

%%%% Das legendäre Schacht-Makro

\usepackage{scalerel}
\newcommand{\mrhup}{\scaleobj{0.5}{\rightharpoonup}}
\newcommand{\mlhup}{\scaleobj{0.5}{\leftharpoonup}}
\newcommand{\mrelb}{\hstretch{0.5}{\scaleobj{0.5}{\relbar}}}
\makeatletter
\def\rhfill@#1{$\m@th\thickmuskip0mu\medmuskip\thickmuskip\thinmuskip\thickmuskip
	\relax#1\mkern+2.5mu\cleaders\hbox{$#1\mrelb\mkern-1mu$}\hfill\mkern-6mu\mrhup\mkern+0.25mu$}
\def\lhfill@#1{$\m@th\thickmuskip0mu\medmuskip\thickmuskip\thinmuskip\thickmuskip
	\relax#1\mkern+1.5mu\mlhup \cleaders\hbox{$#1\mkern-2.5mu\mrelb$}\hfil\mkern+0.5mu$}
\DeclareRobustCommand{\overrightharpoon}{\mathpalette{\overarrow@\rhfill@}}
\DeclareRobustCommand{\overleftharpoon}{\mathpalette{\overarrow@\lhfill@}}
\makeatother

\let\vect=\overrightharpoon

%%% Weiterer Kram von CR
\let\sm=\smallsetminus
\newcommand{\ff}{\mathfrak{f}}

\def\rmlabel{\upshape({\itshape \roman*\,})}
\def\RMlabel{\upshape(\Roman*)}
\def\alabel{\upshape({\itshape \alph*\,})}
\def\Alabel{\upshape({\itshape \Alph*\,})}
\def\nlabel{\upshape({\itshape \arabic*\,})}

\newcommand{\ter}{{\operatorname{ter}}}
\newcommand{\init}{{\operatorname{init}}}
\newcommand{\Idd}{{\operatorname{Id}}}

\let\prt=\partimes
\newcommand{\fJ}{\mathfrak{J}}
\newcommand{\Rho}{\mathrm{P}}

%%%% CUSTOMIZATION 
%%% Margins
\geometry{marginparwidth=21mm, % + <- Width of your marginpar
marginparsep=1mm, left=24mm,right=24mm, top=24mm, bottom=24mm}

%%%%%% Adds hyperlinks
\usepackage[hyperfootnotes=false,draft=false,pdftex, colorlinks=true, linkcolor=black, citecolor=black, urlcolor=black, hypertexnames=false
% pagebackref,
%bookmarksnumbered=true
]{hyperref}

%%% --- 
%%% Catex Package (https://plmlab.math.cnrs.fr/guiraud/catex)
%%% --- 

%%%%%%%%
%% XYPIC %%
%%%%%%%%

%% Options 
\UseTips
\SelectTips{eu}{11}

%% Nouvelle direction pour les monomorphismes (>->)
\newdir{ >}{{}*!/-10pt/@{>}}
\newdir{ -}{{}*!/-10pt/@{}}
\newdir{> }{{}*!/+10pt/@{>}}

%% Ajouts de la variante 4 pour les blancs, les lignes, les pointillés, 
%% les zigouigouis (pas top) et les têtes de flèches (presque finies).
%% Concrètement, dans \xymatrix, la commande \ar @4 produit une 4-cellule.
\makeatletter

% Noms des nouvelles directions
\xyletcsnamecsname@{dir4{}}{dir{}}
\xydefcsname@{dir4{-}}{\line@ \quadruple@\xydashh@}
\xydefcsname@{dir4{.}}{\point@ \quadruple@\xydashh@}
\xydefcsname@{dir4{~}}{\squiggle@ \quadruple@\xybsqlh@}
\xydefcsname@{dir4{>}}{\Tttip@}
\xydefcsname@{dir4{<}}{\reverseDirection@\Tttip@}

% Commande quadruple
% Attention : la position des pourcents est importante !!
\xydef@\quadruple@#1{%
	\edef\Drop@@{%
		\dimen@=#1\relax
		\dimen@=.5\dimen@
		\A@=-\sinDirection\dimen@
		\B@=\cosDirection\dimen@
		\setboxz@h{%
			\setbox2=\hbox{\kern3\A@\raise3\B@\copy\z@}%
			\dp2=\z@ \ht2=\z@ \wd2=\z@ \box2
			\setbox2=\hbox{\kern\A@\raise\B@\copy\z@}%
			\dp2=\z@ \ht2=\z@ \wd2=\z@ \box2
			\setbox2=\hbox{\kern-\A@\raise-\B@\copy\z@}%
			\dp2=\z@ \ht2=\z@ \wd2=\z@ \box2
			\setbox2=\hbox{\kern-3\A@\raise-3\B@ \noexpand\boxz@}%
			\dp2=\z@ \ht2=\z@ \wd2=\z@ \box2
		}%
		\ht\z@=\z@ \dp\z@=\z@ \wd\z@=\z@ \noexpand\styledboxz@
	}%
}

% Têtes de quadruples flèches : positionnement ok, reste la jointure
% Les parties commentées sont des essais à reprendre plus tard
\xydef@\Tttip@{\kern2pt \vrule height2pt depth2pt width\z@
	\Tttip@@ \kern2pt \egroup
	\U@c=0pt \D@c=0pt \L@c=0pt \R@c=0pt \Edge@c={\circleEdge}%
	\def\Leftness@{.5}\def\Upness@{.5}%
	\def\Drop@@{\styledboxz@}\def\Connect@@{\straight@{\dottedSpread@\jot}}}

\xydef@\Tttip@@{%
	\dimen@=.25\dimen@
	%	\A@=-\sinDirection\dimen@
	\B@=\cosDirection\dimen@
	\setboxz@h\bgroup\reverseDirection@\line@ \wdz@=\z@ \ht\z@=\z@ \dp\z@=\z@
	%	\kern\A@ \raise\B@ \boxz@ \kern\L@c
	%	\kern-\L@c \boxz@ \kern\L@c
	{\vDirection@(1,-1)\xydashl@ \xyatipfont\char\DirectionChar}%
	{\vDirection@(1,+1)\xydashl@ \xybtipfont\char\DirectionChar}%
}

% Redéfinition de la commande \ar 
% Provoque un avertissement à la compilation
\xydef@\ar@form{
	\ifx \space@\next \expandafter\DN@\space{\xyFN@\ar@form}%
	\else\ifx ^\next \DN@ ^{\xyFN@\ar@style}\edef\arvariant@@{\string^}%
	\else\ifx _\next \DN@ _{\xyFN@\ar@style}\edef\arvariant@@{\string_}%
	\else\ifx 0\next \DN@ 0{\xyFN@\ar@style}\def\arvariant@@{0}%
	\else\ifx 1\next \DN@ 1{\xyFN@\ar@style}\def\arvariant@@{1}%
	\else\ifx 2\next \DN@ 2{\xyFN@\ar@style}\def\arvariant@@{2}%
	\else\ifx 3\next \DN@ 3{\xyFN@\ar@style}\def\arvariant@@{3}%
	\else\ifx 4\next \DN@ 4{\xyFN@\ar@style}\def\arvariant@@{4}%
	\else\ifx \bgroup\next \let\next@=\ar@style
	\else\ifx [\next \DN@[##1]{\ar@modifiers{[##1]}}%]
	\else\ifx *\next \DN@ *{\ar@modifiers}%
	\else\addLT@\ifx\next \let\next@=\ar@slide
	\else\ifx /\next \let\next@=\ar@curveslash
	\else\ifx (\next \let\next@=\ar@curveinout %)
	\else\addRQ@\ifx\next \addRQ@\DN@{\ar@curve@}%
	\else\addLQ@\ifx\next \addLQ@\DN@{\xyFN@\ar@curve}%
	\else\addDASH@\ifx\next \addDASH@\DN@{\defarstem@-\xyFN@\ar@}%
	\else\addEQ@\ifx\next \addEQ@\DN@{\def\arvariant@@{2}\defarstem@-\xyFN@\ar@}%
	\else\addDOT@\ifx\next \addDOT@\DN@{\defarstem@.\xyFN@\ar@}%
	\else\ifx :\next \DN@:{\def\arvariant@@{2}\defarstem@.\xyFN@\ar@}%
	\else\ifx ~\next \DN@~{\defarstem@~\xyFN@\ar@}%
	\else\ifx !\next \DN@!{\dasharstem@\xyFN@\ar@}%
	\else\ifx ?\next \DN@?{\ar@upsidedown\xyFN@\ar@}%
	\else \let\next@=\ar@error
	\fi\fi\fi\fi\fi\fi\fi\fi\fi\fi\fi\fi\fi\fi\fi\fi\fi\fi\fi\fi\fi\fi\fi \next@}

\makeatother

%LD coherence diagrams

%T1
\newcommand{\Tone}
{\xymatrix@R=0.3cm@C=0cm{
	1\ot (A \prt B) \ar[rr]^{\distl} \ar[ddr]_{\lambda} && (1\ot A)\prt B \ar[ddl]^{\lambda \partimes \id} \\
	& {\lower -5pt\hbox{(T1)}} & \\
	& A\prt B & \\
}\label{eq:T1}
}

%T2
\newcommand{\Ttwo}
{\xymatrix@R=0.3cm@C=0cm{
		(A \prt B) \ot 1 \ar[rr]^{\distr} \ar[ddr]_{\ro} && A \prt {(B\otimes 1)} \ar[ddl]^{\id \prt \ro} \\
		& {\lower -5pt\hbox{(T2)}} & \\
		& A\prt B & \\
	}\label{eq:T2} 
}

%T3
\newcommand{\Tthree}
{\xymatrix@R=0.3cm@C=0cm{
		(K \prt A) \ot B \ar[rr]^{\distr} \ar[ddr]_{\lp \,\ot\, \id} && K \prt {(A \ot B)} \ar[ddl]^{\lp} \\
		& {\lower -5pt\hbox{(T3)}} & \\
		& A\ot B & \\
	}\label{eq:T3} 
}

%T4
\newcommand{\Tfour}
{\xymatrix@R=0.3cm@C=0cm{
		A \ot (B \prt K) \ar[rr]^{\distl} \ar[ddr]_{\id \,\ot\, \rp} && {(A\ot B)} \prt K \ar[ddl]^{\rp} \\
		& {\lower -5pt\hbox{(T4)}} & \\
		& A\ot B & \\
	}\label{eq:T4} 
}

% Catex commands
%% Cellules Catex
\deftwocell[black]{ot : 2 -> 1}
\deftwocell[black,circle]{oteta : 0 -> 1}
\deftwocell[red]{par : 2 -> 1}
\deftwocell[red,circle]{pareta : 0 -> 1}
\deftwocell[crossing]{tau : 2 -> 2}
\deftwocell[gray]{alpha : 3 -> 1}
\deftwocell[gray,lefthalfcircle]{lambda : 1 -> 1}
\deftwocell[gray,righthalfcircle]{rho : 1 -> 1}
\deftwocell[gray,rectangle]{beta : 2 -> 1}
\deftwocell[gray,rectangle]{gamma : 3 -> 1}
\deftwocell[white]{aleph : 4 -> 1}
\deftwocell[white]{beth : 2 -> 1}
\deftwocell[white,rectangle]{gimmel : 2 -> 1}
\deftwocell[white,rectangle]{daleth : 3 -> 1}
\deftwocell[black]{id : 1 -> 1}

\deftwocell[dots]{dots : 2 -> 2}
\deftwocell[gray,rectangle]{f : 2 -> 2}

%% Surface diagrams
\tikzset{anchorbase/.style={baseline={([yshift=-0.5ex]current bounding box.center)}},
	tinynodes/.style={font=\tiny,text height=0.75ex,text depth=0.15ex},
	smallnodes/.style={font=\scriptsize,text height=0.75ex,text depth=0.15ex},
	>={Latex[length=1mm, width=1.5mm]}
}

\newcommand{\aosurf}{
	\adjustbox{margin=0.05cm}{
	\begin{tikzpicture} [anchorbase,xscale=-.2,yscale=.2,tinynodes]
		%bottom web
		\draw [thick] (1,2) to [out=10,in=160] (7,1.5);
		\draw [thick] (2,1) to [out=10,in=150] (5,.5);
		\draw [thick] (3,0) to [out=0,in=210] (5,.5);
		\draw [thick] (5,.5) to [out=0,in=260] (7,1.5);
		\draw [thick] (7,1.5) to (8,1.5);
		%vertical edges
		\draw[thick]  (1,2) to (1,2+4);
		\draw[thick]  (2,1) to (2,1+4);
		\draw[thick]  (3,0) to (3,0+4);
		\draw[thick]  (8,1.5) to (8,1.5+4);
		% top web
		\draw [thick] (1,2+4) to [out=10,in=160] (4,2+4);
		\draw [thick] (2,1+4) to [out=0,in=230] (4,2+4);
		\draw [thick] (3,0+4) to [out=0,in=230] (7,1.5+4);
		\draw [thick] (4,2+4) to[out=0,in=150] (7,1.5+4);
		\draw [thick] (7,1.5+4) to (8,1.5+4);
		%functor lines
		\draw [line width=.6mm,] (5,.5) to [out=80,in=210] (5.75,0+2.85);
		\draw [line width=.6mm,] (7,1.5) to [out=90,in=-10] (5.75,0+2.85);
		\draw [line width=.6mm,] (5.75,0+2.85) to [out=30,in=270] (7,1.5+4);
		\draw [line width=.6mm,] (5.75,0+2.85) to [out=160,in=270] (4,2+4);
		%coherence vertex
		\filldraw[draw=Purple, fill=Purple, ultra thick] (5.75,0+2.85) circle (0.3cm);
	\end{tikzpicture}
}
}

\newcommand{\apsurf}{
\adjustbox{margin=0.05cm}{
	\begin{tikzpicture} [anchorbase,xscale=-.2,yscale=.2,tinynodes]
		%bottom web
		\draw [thick] (1,2) to [out=10,in=160] (7,1.5);
		\draw [thick] (2,1) to [out=10,in=150] (5,.5);
		\draw [thick] (3,0) to [out=0,in=210] (5,.5);
		\draw [thick] (5,.5) to [out=0,in=260] (7,1.5);
		\draw [thick] (7,1.5) to (8,1.5);
		%vertical edges
		\draw[thick]  (1,2) to (1,2+4);
		\draw[thick]  (2,1) to (2,1+4);
		\draw[thick]  (3,0) to (3,0+4);
		\draw[thick]  (8,1.5) to (8,1.5+4);
		% top web
		\draw [thick] (1,2+4) to [out=10,in=160] (4,2+4);
		\draw [thick] (2,1+4) to [out=0,in=230] (4,2+4);
		\draw [thick] (3,0+4) to [out=0,in=230] (7,1.5+4);
		\draw [thick] (4,2+4) to[out=0,in=150] (7,1.5+4);
		\draw [thick] (7,1.5+4) to (8,1.5+4);
		%functor lines
		\draw [line width=.6mm,color=red] (5,.5) to [out=80,in=210] (5.75,0+2.85);
		\draw [line width=.6mm,color=red] (7,1.5) to [out=90,in=-10] (5.75,0+2.85);
		\draw [line width=.6mm,color=red] (5.75,0+2.85) to [out=30,in=270] (7,1.5+4);
		\draw [line width=.6mm,color=red] (5.75,0+2.85) to [out=160,in=270] (4,2+4);
		%coherence vertex
		\filldraw[draw=Purple, fill=Purple, ultra thick] (5.75,0+2.85) circle (0.3cm);
	\end{tikzpicture}
}
}

\newcommand{\distlsurf}{
	\adjustbox{margin=0.05cm}{
		\begin{tikzpicture} [anchorbase,xscale=-.2,yscale=.2,tinynodes]
			%bottom web
			\draw [thick] (1,2) to [out=10,in=160] (7,1.5);
			\draw [thick] (2,1) to [out=10,in=150] (5,.5);
			\draw [thick] (3,0) to [out=0,in=210] (5,.5);
			\draw [thick] (5,.5) to [out=0,in=260] (7,1.5);
			\draw [thick] (7,1.5) to (8,1.5);
			%vertical edges
			\draw[thick]  (1,2) to (1,2+4);
			\draw[thick]  (2,1) to (2,1+4);
			\draw[thick]  (3,0) to (3,0+4);
			\draw[thick]  (8,1.5) to (8,1.5+4);
			% top web
			\draw [thick] (1,2+4) to [out=10,in=160] (4,2+4);
			\draw [thick] (2,1+4) to [out=0,in=230] (4,2+4);
			\draw [thick] (3,0+4) to [out=0,in=230] (7,1.5+4);
			\draw [thick] (4,2+4) to[out=0,in=150] (7,1.5+4);
			\draw [thick] (7,1.5+4) to (8,1.5+4);
			%functor lines
			\draw [line width=.6mm,color=black] (5,.5) to [out=80,in=210] (5.75,0+2.85);
			\draw [line width=.6mm,color=red] (7,1.5) to [out=90,in=-10] (5.75,0+2.85);
			\draw [line width=.6mm,color=black] (5.75,0+2.85) to [out=30,in=270] (7,1.5+4);
			\draw [line width=.6mm,color=red] (5.75,0+2.85) to [out=160,in=270] (4,2+4);
			%coherence vertex
			\filldraw[draw=Purple, fill=Purple, ultra thick] (5.75,0+2.85) circle (0.3cm);
		\end{tikzpicture}
	}
}

\newcommand{\distlsurff}{
\adjustbox{margin=0.05cm}{
	\begin{tikzpicture} [anchorbase,xscale=-.2,yscale=.2,tinynodes]
		%bottom web
		\draw [thick] (1,2) to [out=10,in=160] (4,2);
		\draw [thick] (2,1) to [out=0,in=230] (4,2);
		\draw [thick] (3,0) to [out=0,in=230] (7,1.5);
		\draw [thick] (4,2) to[out=0,in=150] (7,1.5);
		\draw [thick] (7,1.5) to (8,1.5);
		%vertical edges
		\draw[thick]  (1,2) to (1,2+4);
		\draw[thick]  (2,1) to (2,1+4);
		\draw[thick]  (3,0) to (3,0+4);
		\draw[thick]  (8,1.5) to (8,1.5+4);
		% top web
		\draw [thick] (1,2+4) to [out=10,in=160] (7,1.5+4);
		\draw [thick] (2,1+4) to [out=10,in=150] (5,.5+4);
		\draw [thick] (3,0+4) to [out=0,in=210] (5,.5+4);
		\draw [thick] (5,.5+4) to [out=0,in=260] (7,1.5+4);
		\draw [thick] (7,1.5+4) to (8,1.5+4);
		%functor lines
		\draw [line width=.6mm,color=red] (4,2) to [out=90,in=210] (5.75,0+3);
		\draw [line width=.6mm,color=black] (7,1.5) to [out=90,in=-10] (5.75,0+3);
		\draw [line width=.6mm,color=red] (5.75,0+3) to [out=30,in=270] (7,1.5+4);
		\draw [line width=.6mm,color=black] (5.75,0+3) to [out=160,in=270] (5,.5+4);
		%coherence vertex
		\filldraw[draw=Purple, fill=Purple, ultra thick] (5.75,0+3) circle (0.3cm);
	\end{tikzpicture}
}
}

\newcommand{\distrsurf}{
\adjustbox{margin=0.05cm}{
	\begin{tikzpicture} [anchorbase,xscale=-.2,yscale=.2,tinynodes]
		%bottom web
		\draw [thick] (1,2) to [out=10,in=160] (7,1.5);
		\draw [thick] (2,1) to [out=10,in=150] (5,.5);
		\draw [thick] (3,0) to [out=0,in=210] (5,.5);
		\draw [thick] (5,.5) to [out=0,in=260] (7,1.5);
		\draw [thick] (7,1.5) to (8,1.5);
		%vertical edges
		\draw[thick]  (1,2) to (1,2+4);
		\draw[thick]  (2,1) to (2,1+4);
		\draw[thick]  (3,0) to (3,0+4);
		\draw[thick]  (8,1.5) to (8,1.5+4);
		% top web
		\draw [thick] (1,2+4) to [out=10,in=160] (4,2+4);
		\draw [thick] (2,1+4) to [out=0,in=230] (4,2+4);
		\draw [thick] (3,0+4) to [out=0,in=230] (7,1.5+4);
		\draw [thick] (4,2+4) to[out=0,in=150] (7,1.5+4);
		\draw [thick] (7,1.5+4) to (8,1.5+4);
		%functor lines
		\draw [line width=.6mm,color=red] (5,.5) to [out=80,in=210] (5.75,0+2.85);
		\draw [line width=.6mm,color=black] (7,1.5) to [out=90,in=-10] (5.75,0+2.85);
		\draw [line width=.6mm,color=red] (5.75,0+2.85) to [out=30,in=270] (7,1.5+4);
		\draw [line width=.6mm,color=black] (5.75,0+2.85) to [out=160,in=270] (4,2+4);
		%coherence vertex
		\filldraw[draw=Purple, fill=Purple, ultra thick] (5.75,0+2.85) circle (0.3cm);
	\end{tikzpicture}
}
}

\newcommand{\gammalsurf}{
	\adjustbox{margin=0.05cm}{
		\begin{tikzpicture} [anchorbase,xscale=-.2,yscale=.2,tinynodes]
			%bottom web
			\draw [thick] (1,2) to [out=10,in=160] (4,1.5);
			\draw [thick] (4,1.5) to [out=10,in=160] (7,1.5);
			\draw [thick] (2,1) to [out=-10,in=-130] (4,1.5);
			\draw [thick] (3,0) to [out=0,in=240] (7,1.5);
			\draw [thick] (7,1.5) to (8,1.5);
			% top web
			\draw [thick] (1,2+4) to [out=10,in=160] (7,1.5+4);
			\draw [thick] (3,0+4) to [out=0,in=230] (7,1.5+4);
			\draw [thick] (7,1.5+4) to (8,1.5+4);
			%functor lines
			\draw [line width=.6mm,color=red] (4,1.5) to [out=80,in=210] (5.75,0+3.3);
			\draw [line width=.6mm,color=black] (7,1.5) to [out=90,in=-10] (5.75,0+3.3);
			\draw [line width=.6mm,color=red] (5.75,0+3.3) to [out=30,in=270] (7,1.5+4);
			\draw [line width=.6mm,color=black] (5.75,0+3.3) to [out=170,in=90] (2,1);
			%coherence vertex
			\filldraw[draw=Purple, fill=Purple, ultra thick] (5.75,0+3.3) circle (0.3cm);
			%vertical edges
			\draw[thick]  (1,2) to (1,2+4);
			\draw[thick]  (3,0) to (3,0+4);
			\draw[thick]  (8,1.5) to (8,1.5+4);
		\end{tikzpicture}
	}
}

\newcommand{\gammarsurf}{
	\adjustbox{margin=0.05cm}{
		\begin{tikzpicture} [anchorbase,xscale=-.2,yscale=.2,tinynodes]
			%bottom web
			\draw [thick] (1,2) to [out=10,in=160] (4,1.5);
			\draw [thick] (4,1.5) to [out=10,in=160] (7,1.5);
			\draw [thick] (2,1) to [out=-10,in=-130] (4,1.5);
			\draw [thick] (3,0) to [out=0,in=240] (7,1.5);
			\draw [thick] (7,1.5) to (8,1.5);
			% top web
			\draw [thick] (1,2+4) to [out=10,in=160] (7,1.5+4);
			\draw [thick] (3,0+4) to [out=0,in=230] (7,1.5+4);
			\draw [thick] (7,1.5+4) to (8,1.5+4);
			%functor lines
			\draw [line width=.6mm,color=black] (4,1.5) to [out=80,in=210] (5.75,0+3.3);
			\draw [line width=.6mm,color=red] (7,1.5) to [out=90,in=-10] (5.75,0+3.3);
			\draw [line width=.6mm,color=black] (5.75,0+3.3) to [out=30,in=270] (7,1.5+4);
			\draw [line width=.6mm,color=red] (5.75,0+3.3) to [out=170,in=90] (2,1);
			%coherence vertex
			\filldraw[draw=Purple, fill=Purple, ultra thick] (5.75,0+3.3) circle (0.3cm);
			%vertical edges
			\draw[thick]  (1,2) to (1,2+4);
			\draw[thick]  (3,0) to (3,0+4);
			\draw[thick]  (8,1.5) to (8,1.5+4);
		\end{tikzpicture}
	}
}

\newcommand{\losurf}{
\adjustbox{margin=0.05cm}{
	\begin{tikzpicture} [anchorbase,xscale=-.2,yscale=.2,tinynodes]
		%bottom web
		\draw [thick] (1,2) to [out=10,in=160] (6,1.5);
		\draw [thick] (4,.5) to [out=0,in=260] (6,1.5);
		\draw [thick] (6,1.5) to (8,1.65);
		%vertical edges
		\draw[thick]  (1,2) to (1,2+3);
		\draw[thick]  (8,1.65) to (8,1.5+4);
		% top web
		\draw [thick] (1,1+4) to (8,1.5+4);
		%functor lines
		\draw [line width=.6mm,color=black] (4,.5) to [out=90,in=210] (4.75,0+3.5);
		\draw [line width=.6mm,color=black] (6,1.5) to [out=90,in=-10] (4.75,0+3.5);
		%coherence vertex
		\filldraw[draw=Purple, fill=Purple, ultra thick] (4.75,0+3.5) circle (0.3cm);
	\end{tikzpicture}
}
}

\newcommand{\rosurf}{
	\adjustbox{margin=0.05cm}{
	\begin{tikzpicture} [anchorbase,xscale=-.2,yscale=.2,tinynodes]
		%bottom web
		\draw [thick] (1,0.25) to [out=0,in=230] (6,1.5);
		\draw [thick] (3,2) to [out=30,in=110] (6,1.5);
		\draw [thick] (6,1.5) to (8,1.65);
		%vertical edges
		\draw[thick]  (1,0.25) to (1,2+3);
		\draw[thick]  (8,1.65) to (8,1.5+4);
		% top web
		\draw [thick] (1,1+4) to (8,1.5+4);
		%functor lines
		\draw [line width=.6mm,color=black] (3,2) to [out=90,in=180] (4.75,0+3.75);
		\draw [line width=.6mm,color=black] (6,1.5) to [out=90,in=-10] (4.75,0+3.75);
		%coherence vertex
		\filldraw[draw=Purple, fill=Purple, ultra thick] (4.75,0+3.75) circle (0.3cm);
	\end{tikzpicture}
}
}

\newcommand{\lpsurf}{
\adjustbox{margin=0.05cm}{
	\begin{tikzpicture} [anchorbase,xscale=-.2,yscale=.2,tinynodes]
		%bottom web
		\draw [thick] (1,2) to [out=10,in=160] (6,1.5);
		\draw [thick] (4,.5) to [out=0,in=260] (6,1.5);
		\draw [thick] (6,1.5) to (8,1.65);
		%vertical edges
		\draw[thick]  (1,2) to (1,2+3);
		\draw[thick]  (8,1.65) to (8,1.5+4);
		% top web
		\draw [thick] (1,1+4) to (8,1.5+4);
		%functor lines
		\draw [line width=.6mm,color=red] (4,.5) to [out=90,in=210] (4.75,0+3.5);
		\draw [line width=.6mm,color=red] (6,1.5) to [out=90,in=-10] (4.75,0+3.5);
		%coherence vertex
		\filldraw[draw=Purple, fill=Purple, ultra thick] (4.75,0+3.5) circle (0.3cm);
	\end{tikzpicture}
}
}

\newcommand{\rpsurf}{
\adjustbox{margin=0.05cm}{
	\begin{tikzpicture} [anchorbase,xscale=-.2,yscale=.2,tinynodes]
		%bottom web
		\draw [thick] (1,0.25) to [out=0,in=230] (6,1.5);
		\draw [thick] (3,2) to [out=30,in=110] (6,1.5);
		\draw [thick] (6,1.5) to (8,1.65);
		%vertical edges
		\draw[thick]  (1,0.25) to (1,2+3);
		\draw[thick]  (8,1.65) to (8,1.5+4);
		% top web
		\draw [thick] (1,1+4) to (8,1.5+4);
		%functor lines
		\draw [line width=.6mm,color=red] (3,2) to [out=90,in=180] (4.75,0+3.75);
		\draw [line width=.6mm,color=red] (6,1.5) to [out=90,in=-10] (4.75,0+3.75);
		%coherence vertex
		\filldraw[draw=Purple, fill=Purple, ultra thick] (4.75,0+3.75) circle (0.3cm);
	\end{tikzpicture}
}
}

%Whisker
\newcommand{\whisker}{
	\scalebox{1.2}{
	\begin{tikzpicture} [anchorbase,xscale=-.2,yscale=.2,smallnodes]

		%cube f1
		%\draw [thick,lightgray] (17.5,3.9) to (12.5,2.6);
		\draw [thick,lightgray] (17.5,3.95) to (17.5,-0.55);
		\draw [thick,lightgray] (17.5,-0.5) to (12.45,-1.9);
		%\draw [thick,lightgray] (12.5,-1.9) to (12.5,2.6);
		
		\draw [thick,lightgray] (14.53,5.5) to (9.65,3.95);
		\draw [thick,lightgray] (14.5,5.5) to (14.5,0.8);
		\draw [thick,lightgray] (14.5,0.8) to (9.7,-0.6);
		\draw[thick,lightgray]  (9.7,-0.65) to (9.7,4);
		
		\draw[thick,lightgray]  (9.7,-0.6) to (12.5,-1.9);
		\draw[thick,lightgray]  (9.7,4) to (12.5,2.6);
		
		\draw[thick,lightgray]  (17.5,3.9) to (14.5,5.5);
		\draw [thick,lightgray] (17.5,-0.5) to (14.5,0.8);
		\node[lightgray] at (16, 2){$v_1$};
	
	%cube g1

\draw [thick,lightgray] (2.5,11.5) to (-2.5,10.1);
\draw [thick,lightgray] (2.5,11.56) to (2.5,7);
\draw [thick,lightgray] (2.5,7) to (-2.5,5.6);
\draw [thick,lightgray] (-2.5,5.55) to (-2.5,10.15);


%\draw [thick,lightgray] (5.5,9.9) to (0.7,8.4);
\draw [thick,lightgray] (5.5,9.95) to (5.5,5.15);
\draw [thick,lightgray] (5.5,5.2) to (0.7,3.8);
%\draw [thick,lightgray] (0.7,3.8) to (0.7,8.4);


\draw [thick,lightgray] (0.7,3.8) to (-2.5,5.6);
\draw [thick,lightgray] (0.7,8.4) to (-2.5,10.1);
%\begin{pgfonlayer}{foreground}
\draw [thick,lightgray] (5.5,9.9) to (2.5,11.5);
%\end{pgfonlayer}{foreground}
\draw [thick,lightgray] (5.5,5.2) to (2.5,7);

\node[lightgray] at (4.1,8){$w_1$};

%f1
	\draw [thick,lightgray] (12.5,-1.9) to (12.5,2.6);
	\draw [thick,lightgray] (17.5,3.9) to (12.5,2.6);
		
%g1
	\draw [thick,lightgray] (0.7,3.75) to (0.7,8.4);
	\draw [thick,lightgray] (5.5,9.9) to (0.7,8.4);
	
		%cubeg2
		\node[lightgray] at (-0.9, -2.6){$w_2$};
\draw [thick,lightgray] (10.5,9) to (25.5,1.5);
\draw [thick,lightgray] (10.5,13.5) to (25.5,6);
\draw[thick,lightgray]  (10.5,8.95) to (10.5,13.55);
\draw[thick,lightgray]  (25.5,1.45) to (25.5,6);
\draw [thick,lightgray] (29,7) to (14,14.5);
\draw [thick,lightgray] (14,14.5) to (10.5,13.5);
\draw [thick,lightgray] (29,7) to (25.5,6);
\draw [thick,lightgray] (29,7.05) to (29,2.45);
\draw [thick,lightgray] (29,2.5) to (25.5,1.5);
\draw[thick,lightgray]  (14,10) to (14,14.5);
\draw [thick,lightgray] (10.5,9) to (14,10);
\draw [thick,lightgray] (14,10) to (29,2.5);
		
%cubef2
	\node[lightgray] at (27.6, 4.9){$v_2$};
	\draw [thick,lightgray] (-18,1.5) to (-3,-6);
	\draw [thick,lightgray] (-14.5,2.4) to (0.5,-5);
	\draw [thick,lightgray] (-18,6) to (-3,-1.5);
	\draw [thick,lightgray]  (-18,1.45) to (-18,6.05);
	\draw [thick,lightgray]  (-3,-6.05) to (-3,-1.45);
	\draw [thick,lightgray] (0.5,-0.5) to (-14.5,6.9);
	\draw [thick,lightgray] (-14.5,6.9) to (-18,6);
	\draw [thick,lightgray] (0.5,-0.5) to (-3,-1.5);
	\draw [thick,lightgray] (0.5,-0.45) to (0.5,-5.05);
	\draw [thick,lightgray] (0.5,-5) to (-3,-6);
	\draw [thick,lightgray]  (-14.5,2.35) to (-14.5,6.95);
	\draw [thick,lightgray] (-18,1.5) to (-14.5,2.4);
	
	%bottom web
	\draw [thick] (4.1,2.2) to [out=30,in=170] (7.5,3);
	\draw [thick] (5.9,1.8) to [out=0,in=260] (7.5,3);
	\draw [thick] (7.5,3) to (9.3,3.5);
	%vertical edges
	\draw[thick]  (4.1,2.14) to (4.1,6.85);
	\draw[thick]  (9.3,3.45) to (9.3,8.26);
	% top web
	\draw [thick] (4.1,6.8) to (9.3,8.2);
	%functor lines
	\draw [line width=.6mm,color=red] (5.9,1.72) to [out=100,in=220] (6.25,5);
	\draw [line width=.6mm,color=black] (7.5,3) to [out=90,in=-10] (6.25,5);
	%coherence vertex
	\filldraw[draw=Purple, fill=Purple, ultra thick] (6.25,5) circle (0.3cm) node[above=-1, Purple]{$x$};
	
%coordinates
	\draw[thick,->,gray] (0-22-4,0+5+2.4,0-10) -- (4-22-4,1+5+2.4,0-10) node[above=0.2]{$\star_1$};
	\draw[thick,->,gray] (0-26.05,7.38,0-10) -- (0-26.05,7.38+5,0-10) node[right=-1]{$\star_2$};
	\draw[thick,->,gray] (-28,5.5,-15) -- (0-22,0+5.1,0-10) node[right=3]{$\star_0$};

%Rahmen
%\draw [thick] (29,7) to (-3,-1.5);
%\draw [thick] (29,2.5) to (-3,-6);

%\draw [thick] (12.5,-1.9) to (-2.5,5.6);
%\draw [thick] (12.5,2.6) to (-2.5,10.1);

%\draw [thick] (17.5,-0.5) to (2.5,7);
%\draw [thick] (17.5,4) to (2.5,11.5);

%\draw [thick] (14,10) to (-18,1.5);
%\draw [thick] (14,14.5) to (-18,6);
	\end{tikzpicture}
}
}

%%Partial order example

%1
\newcommand{\Partords}{\scalebox{1}{
		\xymatrix{
			& { \twocell{ (oteta *0 1) *1 ot } }
			\\
			{ \twocell{ (oteta *0 1 *0 oteta) *1 (1 *0 ot) *1 ot} }
			\ar@3 [ur] ^-{f_1\eq r \,\star_{1}\, s_2(l)}
			\ar@3 [rr] _-{\vphantom{\rule{0pt}{1.6ex}} f_2\eq s_2(r) \,\star_{1}\, l} ^-{}="1"
			&& { \twocell{ (1 *0 oteta) *1 ot } }
		}
	}
}

%2
\newcommand{\Partordb}{\scalebox{1}{
		\xymatrix{
			& {  \twocell{ (oteta *0 oteta) *1 ot } }
			\\
			{ \twocell{ (oteta *0 oteta *0 oteta) *1 (1 *0 ot) *1 ot} }
			\ar@3 [ur] ^-{g_1 \eq v_2 \,\star_1\, r \,\star_{1}\, s_2(l)}
			\ar@3 [rr] _-{\vphantom{\rule{0pt}{1.6ex}} g_2 \eq v_2 \,\star_1\, s_2(r) \,\star_{1}\, l} ^-{}="1"
			&& { \twocell{ (oteta *0 oteta) *1 ot } }\;.
		}
	}
}

%%Critical branchings

%%Non-confluent critical branchings

%1
\newcommand{\Critdistlronc}{\scalebox{0.92}{
		\xymatrix{
			& { \twocell{ ( 1 *0 oteta *0 1) *1 (1 *0 par) *1 ot } }
			\\
			{ \twocell{ ( 1 *0 oteta *0 1) *1 (ot *0 1) *1 par} }
			\ar@3 [ur] ^-{\distlc}
			\ar@3 [rr] _-{\vphantom{\rule{0pt}{1.2ex}} \roc} ^-{}="1"
			&& { \twocell{par} }\;,
		}
	}
}

%2
\newcommand{\Critdistrrpnc}{\scalebox{0.92}{
		\xymatrix{
			& { \twocell{ ( 1 *0 pareta *0 1) *1 (1 *0 ot) *1 par } }
			\\
			{ \twocell{ ( 1 *0 pareta *0 1) *1 (par *0 1) *1 ot} }
			\ar@3 [ur] ^-{\distrc}
			\ar@3 [rr] _-{\vphantom{\rule{0pt}{1.6ex}} \rpc} ^-{}="1"
			&& { \twocell{ot} }\;.
		}
	}
}

%% Confluent critical branchings of LD
%1
\newcommand{\Critdistllo}{\scalebox{0.9}{
	\xymatrix{
		& { \twocell{ ( oteta *0 1 *0 1) *1 (1 *0 par) *1 ot} }
		\ar@[lightgray]@3[dr] ^-{\color{lightgray}\loc}
		\\
		{ \twocell{ ( oteta *0 1 *0 1) *1 (ot *0 1) *1 par } }
		\ar@3 [ur] ^-{\distlc}
		\ar@3 [rr] _-{ \loc \vphantom{\rule{0pt}{1.6ex}}} ^-{}="1"
		&& { \twocell{par} }
		\ar@0 "2,2"!<20pt,15pt>;"1"!<0pt,0pt> _*+{(1)}
	}\label{Crit1}
}
}

%2
\newcommand{\Critdistrro}{\scalebox{0.9}{
		\xymatrix{
			& { \twocell{ ( 1 *0 1 *0 oteta) *1 (1 *0 ot) *1 par } }
			\ar@[lightgray]@3[dr] ^-{\color{lightgray}\roc}
			\\
			{ \twocell{ ( 1 *0 1 *0 oteta) *1 (par *0 1) *1 ot} }
			\ar@3 [ur] ^-{\distrc}
			\ar@3 [rr] _-{\roc  \vphantom{\rule{0pt}{1.2ex}}} ^-{}="1"
			&& { \twocell{par} }
			\ar@0 "2,2"!<20pt,15pt>;"1"!<0pt,0pt> _*+{(2)}
		}\label{Crit2}
	}
}

%3
\newcommand{\Critdistrlp}{\scalebox{0.9}{
		\xymatrix{
			& { \twocell{ ( pareta *0 1 *0 1) *1 (1 *0 ot) *1 par } }
			\ar@[lightgray]@3[dr] ^-{\color{lightgray}\lpc}
			\\
			{ \twocell{ ( pareta *0 1 *0 1) *1 (par *0 1) *1 ot} }
			\ar@3 [ur] ^-{\distrc}
			\ar@3 [rr] _-{\lpc  \vphantom{\rule{0pt}{1.8ex}}} ^-{}="1"
			&& { \twocell{ot} }
			\ar@0 "2,2"!<20pt,15pt>;"1"!<0pt,0pt> _*+{(3)}
		}\label{Crit3}
	}
}

%4
\newcommand{\Critdistlrp}{\scalebox{0.9}{
		\xymatrix{
			& { \twocell{ ( 1 *0 1 *0 pareta) *1 (1 *0 par) *1 ot } }
			\ar@[lightgray]@3[dr] ^-{\color{lightgray}\rpc}
			\\
			{ \twocell{ ( 1 *0 1 *0 pareta) *1 (ot *0 1) *1 par} }
			\ar@3 [ur] ^-{\distlc}
			\ar@3 [rr] _-{\rpc  \vphantom{\rule{0pt}{1.6ex}}} ^-{}="1"
			&& { \twocell{ot} }
			\ar@0 "2,2"!<20pt,15pt>;"1"!<0pt,0pt> _*+{(4)}
		}\label{Crit4}
	}
}

%5
\newcommand{\Critdistlao}{\scalebox{0.9}{
	\xymatrix@R=10pt@C=15pt{
		&{ \twocell{ (1 *0 ot *0 1) *1 (ot *0 1) *1 par } }
		\ar@[lightgray]@3 [rr] 
		^-{\color{lightgray}\distlc} 
		_-{}="1" 
		&& { \twocell{ (1 *0 ot *0 1) *1 (1 *0 par) *1 ot } }
		\ar@[lightgray]@3 [dr] ^-{\color{lightgray}\distlc}
		\\
		{ \twocell{(ot *0 2) *1 (ot *0 1) *1 par } }
		\ar@3 [ur] ^-{\aoc}
		\ar@3 [drr] _-{\distlc  \vphantom{\rule{0pt}{1.7ex}}}
		&&&& { \twocell{ (1 *0 1 *0 par) *1 (1 *0 ot) *1 ot } }
		\\
		&& { \twocell{ (ot *0 par) *1 ot }  }
		\ar@[lightgray]@3 [urr] _-{\color{lightgray}\aoc}
		\ar@0 "3,2"!<25pt,20pt>;"1"!<0pt,0pt> _-*+{(5)}
	}\label{Crit5}
}
}

%6
\newcommand{\Critaodistl}{\scalebox{0.9}{
		\xymatrix@R=10pt@C=15pt{
		&{ \twocell{ (1 *0 par *0 1) *1 (ot *0 1) *1 ot } }
		\ar@[lightgray]@3 [rr] 
		^-{\color{lightgray}\aoc} 
		_-{}="1" 
		&& { \twocell{ (1 *0 par *0 1) *1 (1 *0 ot) *1 ot } }
		\ar@[lightgray]@3 [dr] ^-{\color{lightgray}\distrc}
		\\
		{ \twocell{(ot *0 2) *1 (par *0 1) *1 ot } }
		\ar@3 [ur] ^-{\distlc}
		\ar@3 [drr] _-{\distrc}
		&&&& { \twocell{ (1 *0 1 *0 ot) *1 (1 *0 par) *1 ot } }
		\\
		&& { \twocell{ (ot *0 ot) *1 par }  }
		\ar@[lightgray]@3 [urr] _-{\color{lightgray}\distlc}
		\ar@0 "3,2"!<25pt,20pt>;"1"!<0pt,0pt> _-*+{(6)}\label{Crit6}
	}
	}
}

%7
\newcommand{\Critdistlap}{\scalebox{0.9}{
		\xymatrix@R=10pt@C=15pt{
			&{ \twocell{(1 *0 par *0 1) *1 (ot *0 1) *1 par } }
			\ar@[lightgray]@3 [rr] 
			^-{\color{lightgray}\distlc} 
			_-{}="1" 
			&& { \twocell{ (1 *0 par *0 1) *1 (1 *0 par) *1 ot } }
			\ar@[lightgray]@3 [dr] ^-{\color{lightgray}\apc}
			\\
			{ \twocell{ (ot *0 2) *1 (par *0 1) *1 par } }
			\ar@3 [ur] ^-{\distlc}
			\ar@3 [drr] _-{\apc}
			&&&& { \twocell{ (2 *0 par) *1 (1 *0 par) *1 ot } }
			\\
			&& { \twocell{ (ot *0 par) *1 par } }
			\ar@[lightgray]@3 [urr] _-{\color{lightgray}\distlc}
			\ar@0 "3,2"!<25pt,20pt>;"1"!<0pt,0pt> _-*+{(7)}
		}\label{Crit7}
	}
}

%8
\newcommand{\Critdistrao}{\scalebox{0.9}{
		\xymatrix@R=10pt@C=15pt{
			&{ \twocell{(1 *0 ot *0 1) *1 (par *0 1) *1 ot } }
			\ar@[lightgray]@3 [rr] 
			^-{\color{lightgray}\distrc} 
			_-{}="1" 
			&& { \twocell{ (1 *0 ot *0 1) *1 (1 *0 ot) *1 par } }
			\ar@[lightgray]@3 [dr] ^-{\color{lightgray}\aoc}
			\\
			{ \twocell{ (par *0 2) *1 (ot *0 1) *1 ot } }
			\ar@3 [ur] ^-{\distrc}
			\ar@3 [drr] _-{\aoc}
			&&&& { \twocell{ (2 *0 ot) *1 (1 *0 ot) *1 par } }
			\\
			&& { \twocell{ (par *0 ot) *1 ot } }
			\ar@[lightgray]@3 [urr] _-{\color{lightgray}\distrc}
			\ar@0 "3,2"!<25pt,20pt>;"1"!<0pt,0pt> _-*+{(8)}
		}\label{Crit8}
	}
}

%9
\newcommand{\Critdistldistr}{\scalebox{0.9}{
		\xymatrix@R=10pt@C=15pt{
			&{ \twocell{(1 *0 ot *0 1) *1 (par *0 1) *1 par } }
			\ar@[lightgray]@3 [rr] 
			^-{\color{lightgray}\apc} 
			_-{}="1" 
			&& { \twocell{ (1 *0 ot *0 1) *1 (1 *0 par) *1 par } }
			\ar@[lightgray]@3 [dr] ^-{\color{lightgray}\distrc}
			\\
			{ \twocell{(par *0 2) *1 (ot *0 1) *1 par } }
			\ar@3 [ur] ^-{\distrc}
			\ar@3 [drr] _-{\distlc}
			&&&& {  \twocell{ (2 *0 par) *1 (1 *0 ot) *1 par }  }
			\\
			&& {  \twocell{ (par *0 par) *1 ot } }
			\ar@[lightgray]@3 [urr] _-{\color{lightgray}\distrc}
			\ar@0 "3,2"!<25pt,20pt>;"1"!<0pt,0pt> _-*+{(9)}
		}\label{Crit9}
	}
}

%10
\newcommand{\Critapdistr}{\scalebox{0.9}{
		\xymatrix@R=10pt@C=15pt{
			&{ \twocell{(1 *0 par *0 1) *1 (par *0 1) *1 ot }}
			\ar@[lightgray]@3 [rr] 
			^-{\color{lightgray}\distrc} 
			_-{}="1" 
			&& { \twocell{ (1 *0 par *0 1) *1 (1 *0 ot) *1 par } }
			\ar@[lightgray]@3 [dr] ^-{\color{lightgray}\distrc}
			\\
			{ \twocell{(par *0 2) *1 (par *0 1) *1 ot } }
			\ar@3 [ur] ^-{\apc}
			\ar@3 [drr] _-{\distrc}
			&&&& { \twocell{ (2 *0 ot) *1 (1 *0 par) *1 par } }
			\\
			&& { \twocell{ (par *0 ot) *1 par } }
			\ar@[lightgray]@3 [urr] _-{\color{lightgray}\apc}
			\ar@0 "3,2"!<20pt,20pt>;"1"!<0pt,0pt> _-*+{(10)}
		}\label{Crit10}
	}
}

%15
\newcommand{\Critaoao}{\scalebox{0.9}{
		\xymatrix@R=10pt@C=15pt{
			&{ \twocell{(1 *0 ot *0 1) *1 (ot *0 1) *1 ot }}
			\ar@[lightgray]@3 [rr] 
			^-{\color{lightgray}\aoc} 
			_-{}="1" 
			&& { \twocell{ (1 *0 ot *0 1) *1 (1 *0 ot) *1 ot } }
			\ar@[lightgray]@3 [dr] ^-{\color{lightgray}\aoc}
			\\
			{ \twocell{(ot *0 2) *1 (ot *0 1) *1 ot } }
			\ar@3 [ur] ^-{\aoc}
			\ar@3 [drr] _-{\aoc}
			&&&& { \twocell{ (2 *0 ot) *1 (1 *0 ot) *1 ot } }
			\\
			&& { \twocell{ (ot *0 ot) *1 ot } }
			\ar@[lightgray]@3 [urr] _-{\color{lightgray}\aoc}
			\ar@0 "3,2"!<20pt,20pt>;"1"!<0pt,0pt> _-*+{(15)}
		}\label{Crit15}
	}
}

%11
\newcommand{\Critaoro}{\scalebox{0.9}{
		\xymatrix{
			& { \twocell{ ( 1 *0 oteta *0 1) *1 (1 *0 ot) *1 ot } }
			\ar@[lightgray]@3[dr] ^-{\color{lightgray}\loc}
			\\
			{ \twocell{ ( 1 *0 oteta *0 1) *1 (ot *0 1) *1 ot} }
			\ar@3 [ur] ^-{\aoc}
			\ar@3 [rr] _-{\vphantom{\rule{0pt}{1.6ex}} \roc} ^-{}="1"
			&& { \twocell{ot} }
			\ar@0 "2,2"!<22pt,15pt>;"1"!<0pt,0pt> _*+{(11)}
		}\label{Crit11}
	}
}

%12
\newcommand{\Critaolo}{\scalebox{0.9}{
		\xymatrix{
			& { \twocell{ ( oteta *0 1 *0 1) *1 (1 *0 ot) *1 ot } }
			\ar@[lightgray]@3[dr] ^-{\color{lightgray}\loc}
			\\
			{ \twocell{ ( oteta *0 1 *0 1) *1 (ot *0 1) *1 ot} }
			\ar@3 [ur] ^-{\aoc}
			\ar@3 [rr] _-{\vphantom{\rule{0pt}{1.6ex}} \loc} ^-{}="1"
			&& { \twocell{ot} }
			\ar@0 "2,2"!<22pt,15pt>;"1"!<0pt,0pt> _*+{(12)}
		}\label{Crit12}
	}
}

%13
\newcommand{\Critaororo}{\scalebox{0.9}{
		\xymatrix{
			& { \twocell{ ( 1 *0 1 *0 oteta) *1 (1 *0 ot) *1 ot } }
			\ar@[lightgray]@3[dr] ^-{\color{lightgray}\roc}
			\\
			{ \twocell{ ( 1 *0 1 *0 oteta) *1 (ot *0 1) *1 ot} }
			\ar@3 [ur] ^-{\aoc}
			\ar@3 [rr] _-{\vphantom{\rule{0pt}{1.6ex}} \roc} ^-{}="1"
			&& { \twocell{ot} }
			\ar@0 "2,2"!<22pt,15pt>;"1"!<0pt,0pt> _*+{(13)}
		}\label{Crit13}
	}
}

%14
\newcommand{\Critloro}{\scalebox{0.9}{
	\xymatrix{
		{\twocell{(oteta *0 oteta) *1 ot}}
		\ar@3@/^5ex/ [rr] ^{\vphantom{\rule{0pt}{1.2ex}} \loc} _{}="1"
		\ar@3@/_5ex/ [rr] _{\vphantom{\rule{0pt}{1.2ex}} \roc} ^{}="2"
		&& {\twocell{oteta}}
		\ar@0 "1,1"!<-4pt,-7pt>;"1"!<20pt,15pt> _-*+{(14)}
	}\label{Crit14}
}
}

%20
\newcommand{\Critapap}{\scalebox{0.9}{
		\xymatrix@R=10pt@C=15pt{
			&{ \twocell{(1 *0 par *0 1) *1 (par *0 1) *1 par }}
			\ar@[lightgray]@3 [rr] 
			^-{\color{lightgray}\apc} 
			_-{}="1" 
			&& { \twocell{ (1 *0 par *0 1) *1 (1 *0 par) *1 par } }
			\ar@[lightgray]@3 [dr] ^-{\color{lightgray}\apc}
			\\
			{ \twocell{(par *0 2) *1 (par *0 1) *1 par } }
			\ar@3 [ur] ^-{\apc}
			\ar@3 [drr] _-{\vphantom{\rule{0pt}{1.4ex}}\apc}
			&&&& { \twocell{ (2 *0 par) *1 (1 *0 par) *1 par } }
			\\
			&& { \twocell{ (par *0 par) *1 par } }
			\ar@[lightgray]@3 [urr] _-{\color{lightgray}\apc}
			\ar@0 "3,2"!<20pt,20pt>;"1"!<0pt,0pt> _-*+{(20)}
		}\label{Crit20}
	}
}

%16
\newcommand{\Critaprp}{\scalebox{0.9}{
		\xymatrix{
			& { \twocell{ ( 1 *0 pareta *0 1) *1 (1 *0 par) *1 par } }
			\ar@[lightgray]@3[dr] ^-{\color{lightgray}\lpc}
			\\
			{ \twocell{ ( 1 *0 pareta *0 1) *1 (par *0 1) *1 par} }
			\ar@3 [ur] ^-{\apc}
			\ar@3 [rr] _-{\vphantom{\rule{0pt}{1.4ex}} \rpc} ^-{}="1"
			&& { \twocell{par} }
			\ar@0 "2,2"!<20pt,15pt>;"1"!<0pt,0pt> _*+{(16)}
		}\label{Crit16}
	}
}

%17
\newcommand{\Critaplp}{\scalebox{0.9}{
		\xymatrix{
			& { \twocell{ ( pareta *0 1 *0 1) *1 (1 *0 par) *1 par } }
			\ar@[lightgray]@3[dr] ^-{\color{lightgray}\lpc}
			\\
			{ \twocell{ ( pareta *0 1 *0 1) *1 (par *0 1) *1 par} }
			\ar@3 [ur] ^-{\apc}
			\ar@3 [rr] _-{\vphantom{\rule{0pt}{1.5ex}} \lpc} ^-{}="1"
			&& { \twocell{par} }
			\ar@0 "2,2"!<20pt,15pt>;"1"!<0pt,0pt> _*+{(17)}
		}\label{Crit17}
	}
}

%18
\newcommand{\Critaprprp}{\scalebox{0.9}{
		\xymatrix{
			& { \twocell{ ( 1 *0 1 *0 pareta) *1 (1 *0 par) *1 par } }
			\ar@[lightgray]@3[dr] ^-{\color{lightgray}\rpc}
			\\
			{ \twocell{ ( 1 *0 1 *0 pareta) *1 (par *0 1) *1 par} }
			\ar@3 [ur] ^-{\apc}
			\ar@3 [rr] _-{\vphantom{\rule{0pt}{1.6ex}} \rpc} ^-{}="1"
			&& { \twocell{par} }
			\ar@0 "2,2"!<20pt,15pt>;"1"!<0pt,0pt> _*+{(18)}
		}\label{Crit18}
	}
}

%19
\newcommand{\Critlprp}{\scalebox{0.9}{
		\xymatrix{
			{\twocell{(pareta *0 pareta) *1 par}}
			\ar@3@/^5ex/ [rr] ^{\lpc} _{}="1"
			\ar@3@/_5ex/ [rr] _{\vphantom{\rule{0pt}{1.5ex}} \rpc} ^{}="2"
			&& {\twocell{pareta}}
			\ar@0 "1,1"!<-4pt,-7pt>;"1"!<20pt,15pt> _-*+{(19)}
		}\label{Crit19}
	}
}

%%Resolved critical branching
%21
\newcommand{\Critdistlroncr}{\scalebox{0.9}{
		\xymatrix{
			& { \twocell{ ( 1 *0 oteta *0 1) *1 (1 *0 par) *1 ot } }
			\ar@[lightgray]@3[dr] ^-{\color{lightgray}k^l}
			\\
			{ \twocell{ ( 1 *0 oteta *0 1) *1 (ot *0 1) *1 par} }
			\ar@3 [ur] ^-{\distlc}
			\ar@3 [rr] _-{\vphantom{\rule{0pt}{1.2ex}} \roc} ^-{}="1"
			&& { \twocell{par} }
			\ar@0 "2,2"!<20pt,15pt>;"1"!<0pt,0pt> _*+{(21)}
		}\label{Crit21}
	}
}

%22
\newcommand{\Critdistrrpncr}{\scalebox{0.9}{
		\xymatrix{
			& { \twocell{ ( 1 *0 pareta *0 1) *1 (1 *0 ot) *1 par } }
			\ar@[lightgray]@3[dr] ^-{\color{lightgray}k^r}
			\\
			{ \twocell{ ( 1 *0 pareta *0 1) *1 (par *0 1) *1 ot} }
			\ar@3 [ur] ^-{\distrc}
			\ar@3 [rr] _-{\vphantom{\rule{0pt}{1.5ex}} \rpc} ^-{}="1"
			&& { \twocell{ot} }
			\ar@0 "2,2"!<20pt,15pt>;"1"!<0pt,0pt> _*+{(22)}
		}\label{Crit22}
	}
}

%%New critical branching

%Involving gammal and not gammar

%23
\newcommand{\Critglrp}{\scalebox{0.9}{
		\xymatrix{
			{ \twocell{ ( 1 *0 oteta *0 pareta) *1 (1 *0 par) *1 ot } }
			\ar@3[rr] ^-{k^l}
			\ar@3 [d] ^-{\hspace{2pt} \rpc}
			&& { \twocell{ (1 *0 pareta) *1 par} }
			\ar@[lightgray]@3 [d] ^-{\hspace{2pt} \color{lightgray}\rpc}
			\\
			{ \twocell{ (1 *0 oteta) *1 ot } }
			\ar@[lightgray]@3 [rr] _-{\vphantom{\rule{0pt}{1.5ex}} \color{lightgray}\roc} ^-{}="1"
			&& { \twocell{id *1 id} }
			\ar@0 "2,2"!<30pt,30pt>;"1"!<0pt,0pt> _*+{(23)}
		}\label{Crit23}
	}
}

%24
\newcommand{\Critgllo}{\scalebox{0.9}{
	\xymatrix{
		{\twocell{ ( oteta *0 oteta *0 1) *1 (1 *0 par) *1 ot }}
		\ar@3@/^5ex/ [rr] ^{k^l} _{}="1"
		\ar@3@/_5ex/ [rr] _{\vphantom{\rule{0pt}{1.5ex}} \loc} ^{}="2"
		&& {\twocell{(oteta *0 1) *1 par }}
		\ar@0 "1,1"!<-1pt,-7pt>;"1"!<20pt,15pt> _-*+{(24)}
	}\label{Crit24}
	}
}

%25
\newcommand{\Critgldistr}{\scalebox{0.9}{
		\xymatrix{
			{ \twocell{ ( 2 *0 oteta *0 1) *1 (par *0 par) *1 ot } }
			\ar@3[rr] ^-{k^l}
			\ar@3 [d] ^-{\hspace{2pt}\distrc}
			&& { \twocell{ (par *0 1) *1 par} }
			\ar@[lightgray]@3 [d] ^-{\hspace{2pt}\color{lightgray}\apc}
			\\
			{ \twocell{(2 *0 oteta *0 1) *1 (2 *0 par) *1 (1 *0 ot) *1 par } }
			\ar@[lightgray]@3 [rr] _-{\vphantom{\rule{0pt}{1.5ex}} \color{lightgray}k^l} ^-{}="1"
			&& { \twocell{ (1 *0 par) *1 par } }
			\ar@0 "2,2"!<30pt,30pt>;"1"!<0pt,0pt> _*+{(25)}
		}\label{Crit25}
	}
}

%26
\newcommand{\Critglao}{\scalebox{0.9}{
		\xymatrix{
			{ \twocell{ ( 2 *0 oteta *0 1) *1 (ot *0 par) *1 ot } }
			\ar@3[rr] ^-{k^l}
			\ar@3 [d] ^-{\hspace{2pt} \aoc}
			&& { \twocell{ (ot *0 1) *1 par} }
			\ar@[lightgray]@3 [d] ^-{\hspace{2pt} \color{lightgray}\distlc}
			\\
			{ \twocell{(2 *0 oteta *0 1) *1 (2 *0 par) *1 (1 *0 ot) *1 ot } }
			\ar@[lightgray]@3 [rr] _-{\vphantom{\rule{0pt}{1.5ex}} \color{lightgray}k^l} ^-{}="1"
			&& { \twocell{ (1 *0 par) *1 ot } }
			\ar@0 "2,2"!<30pt,30pt>;"1"!<0pt,0pt> _*+{(26)}
		}\label{Crit26}
	}
}

%27
\newcommand{\Critgldistl}{\scalebox{0.9}{
		\xymatrix@R=10pt@C=15pt{
			&{ \twocell{(1 *0 oteta *0 2) *1 (1 *0 par *0 1) *1 (1 *0 par) *1 ot }}
			\ar@[lightgray]@3 [rr] 
			^-{\color{lightgray}\apc} 
			_-{}="1" 
			&& { \twocell{ (1 *0 oteta *0 par) *1 (1 *0 par) *1 ot } }
			\ar@[lightgray]@3 [dr] ^-{\color{lightgray}k^l}
			\\
			{ \twocell{ (1 *0 oteta *0 2) *1 (1 *0 par *0 1) *1 (ot *0 1) *1 par } }
			\ar@3 [ur] ^-{\distlc}
			\ar@3 [drr] _-{\vphantom{\rule{0pt}{1.7ex}} k^l}
			&&&& { \twocell{ (1 *0 par) *1 par } }
			\\
			&& { \twocell{ (par *0 1) *1 par } }
			\ar@[lightgray]@3 [urr] _-{\color{lightgray}\apc}
			\ar@0 "3,2"!<20pt,20pt>;"1"!<0pt,0pt> _-*+{(27)}
		}\label{Crit27}
	}
}

%28
\newcommand{\CritglaoP}{\scalebox{0.9}{
		\xymatrix@R=10pt@C=15pt{
			&{ \twocell{(1 *0 oteta *0 2) *1 (1 *0 par *0 1) *1 (1 *0 ot) *1 ot }}
			\ar@[lightgray]@3 [rr] 
			^-{\color{lightgray}\distrc} 
			_-{}="1" 
			&& { \twocell{ (1 *0 oteta *0 ot) *1 (1 *0 par) *1 ot } }
			\ar@[lightgray]@3 [dr] ^-{\color{lightgray}k^l}
			\\
			{ \twocell{ (1 *0 oteta *0 2) *1 (1 *0 par *0 1) *1 (ot *0 1) *1 ot } }
			\ar@3 [ur] ^-{\aoc}
			\ar@3 [drr] _-{\vphantom{\rule{0pt}{1.7ex}} k^l}
			&&&& { \twocell{ (1 *0 ot) *1 par } }
			\\
			&& { \twocell{ (par *0 1) *1 ot } }
			\ar@[lightgray]@3 [urr] _-{\color{lightgray}\distrc}
			\ar@0 "3,2"!<20pt,20pt>;"1"!<0pt,0pt> _-*+{(28)}
		}\label{Crit28}
	}
}

%Involving gammar and not gammal

%29
\newcommand{\Critgrro}{\scalebox{0.9}{
		\xymatrix{
			{ \twocell{ ( 1 *0 pareta *0 oteta) *1 (1 *0 ot) *1 par } }
			\ar@3[rr] ^-{k^r}
			\ar@3 [d] ^-{\hspace{2pt} \roc}
			&& { \twocell{ (1 *0 oteta) *1 ot} }
			\ar@[lightgray]@3 [d] ^-{\hspace{2pt} \color{lightgray}\roc}
			\\
			{ \twocell{ (1 *0 pareta) *1 par } }
			\ar@[lightgray]@3 [rr] _-{\vphantom{\rule{0pt}{1.5ex}} \color{lightgray}\rpc} ^-{}="1"
			&& { \twocell{id *1 id} }
			\ar@0 "2,2"!<27pt,22pt>;"1"!<0pt,0pt>_*+{(29)}
		}\label{Crit29}
	}
}

%30
\newcommand{\Critgrlp}{\scalebox{0.9}{
		\xymatrix{
			{\twocell{ ( pareta *0 pareta *0 1) *1 (1 *0 ot) *1 par }}
			\ar@3@/^5ex/ [rr] ^{k^r} _{}="1"
			\ar@3@/_5ex/ [rr] _{\vphantom{\rule{0pt}{1.5ex}} \lpc} ^{}="2"
			&& {\twocell{(pareta *0 1) *1 ot }}
			\ar@0 "1,1"!<0pt,-7pt>;"1"!<20pt,15pt> _-*+{(30)}
		}\label{Crit30}
	}
}

%31
\newcommand{\Critgrdistl}{\scalebox{0.9}{
		\xymatrix{
			{ \twocell{ ( 2 *0 pareta *0 1) *1 (ot *0 ot) *1 par } }
			\ar@3[rr] ^-{k^r}
			\ar@3 [d] ^-{\hspace{2pt} \distlc}
			&& { \twocell{ (ot *0 1) *1 ot} }
			\ar@[lightgray]@3 [d] ^-{\hspace{2pt} \color{lightgray}\aoc}
			\\
			{ \twocell{(2 *0 pareta *0 1) *1 (2 *0 ot) *1 (1 *0 par) *1 ot } }
			\ar@[lightgray]@3 [rr] _-{\vphantom{\rule{0pt}{1.5ex}} \color{lightgray}k^r} ^-{}="1"
			&& { \twocell{ (1 *0 ot) *1 ot } }
			\ar@0 "2,2"!<30pt,30pt>;"1"!<0pt,0pt>_*+{(31)}
		}\label{Crit31}
	}
}

%32
\newcommand{\Critgrap}{\scalebox{0.9}{
		\xymatrix{
			{ \twocell{ ( 2 *0 pareta *0 1) *1 (par *0 ot) *1 par } }
			\ar@3[rr] ^-{k^r}
			\ar@3 [d] ^-{\hspace{2pt} \apc}
			&& { \twocell{ (par *0 1) *1 ot} }
			\ar@[lightgray]@3 [d] ^-{\hspace{2pt} \color{lightgray}\distrc}
			\\
			{ \twocell{(2 *0 pareta *0 1) *1 (2 *0 ot) *1 (1 *0 par) *1 par } }
			\ar@[lightgray]@3 [rr] _-{\vphantom{\rule{0pt}{1.5ex}} \color{lightgray}k^r} ^-{}="1"
			&& { \twocell{ (1 *0 ot) *1 par } }
			\ar@0 "2,2"!<30pt,30pt>;"1"!<0pt,0pt>_*+{(32)}
		}\label{Crit32}
	}
}

%33
\newcommand{\Critgrdistr}{\scalebox{0.9}{
		\xymatrix@R=10pt@C=15pt{
			&{ \twocell{(1 *0 pareta *0 2) *1 (1 *0 ot *0 1) *1 (1 *0 ot) *1 par }}
			\ar@[lightgray]@3 [rr] 
			^-{\color{lightgray}\aoc} 
			_-{}="1" 
			&& { \twocell{ (1 *0 pareta *0 ot) *1 (1 *0 ot) *1 par } }
			\ar@[lightgray]@3 [dr] ^-{\color{lightgray}k^r}
			\\
			{ \twocell{ (1 *0 pareta *0 2) *1 (1 *0 ot *0 1) *1 (par *0 1) *1 ot } }
			\ar@3 [ur] ^-{\distrc}
			\ar@3 [drr] _-{\vphantom{\rule{0pt}{1.5ex}} k^r}
			&&&& { \twocell{ (1 *0 ot) *1 ot } }
			\\
			&& { \twocell{ (ot *0 1) *1 ot } }
			\ar@[lightgray]@3 [urr] _-{\color{lightgray}\aoc}
			\ar@0 "3,2"!<20pt,20pt>;"1"!<0pt,0pt> _-*+{(33)}
		}\label{Crit33}
	}
}

%34
\newcommand{\CritgrapP}{\scalebox{0.9}{
		\xymatrix@R=10pt@C=15pt{
			&{ \twocell{(1 *0 pareta *0 2) *1 (1 *0 ot *0 1) *1 (1 *0 par) *1 par }}
			\ar@[lightgray]@3 [rr] 
			^-{\color{lightgray}\distlc} 
			_-{}="1" 
			&& { \twocell{ (1 *0 pareta *0 par) *1 (1 *0 ot) *1 par } }
			\ar@[lightgray]@3 [dr] ^-{\color{lightgray}k^r}
			\\
			{ \twocell{ (1 *0 pareta *0 2) *1 (1 *0 ot *0 1) *1 (par *0 1) *1 par } }
			\ar@3 [ur] ^-{\apc}
			\ar@3 [drr] _-{\vphantom{\rule{0pt}{1.5ex}} k^r}
			&&&& { \twocell{ (1 *0 par) *1 ot } }
			\\
			&& { \twocell{ (ot *0 1) *1 par } }
			\ar@[lightgray]@3 [urr] _-{\color{lightgray}\distlc}
			\ar@0 "3,2"!<20pt,20pt>;"1"!<0pt,0pt> _-*+{(34)}
		}\label{Crit34}
	}
}

%Involving both gammal and gammar

%35
\newcommand{\Critglgro}{\scalebox{0.9}{
		\xymatrix{
			{ \twocell{ (1 *0 oteta *0 pareta *0 1) *1 ( 2 *0 ot) *1 (1 *0 par) *1 ot } }
			\ar@3[rr] ^-{k^l}
			\ar@3 [d] ^-{\hspace{2pt} k^r}
			&& { \twocell{ (1 *0 pareta *0 1) *1 (1 *0 ot) *1 par} }
			\ar@[lightgray]@3 [d] ^-{\hspace{2pt} \color{lightgray}k^r}
			\\
			{ \twocell{(1 *0 oteta *0 1) *1 (1 *0 ot) *1 ot } }
			\ar@[lightgray]@3 [rr] _-{\vphantom{\rule{0pt}{1.5ex}} \color{lightgray}\loc} ^-{}="1"
			&& { \twocell{ot} }
			\ar@0 "2,2"!<30pt,30pt>;"1"!<0pt,0pt>_*+{(35)}
		}\label{Crit35}
	}
}

%36
\newcommand{\Critglgrt}{\scalebox{0.9}{
	\xymatrix{
		{ \twocell{ (1 *0 pareta *0 oteta *0 1) *1 ( 2 *0 par) *1 (1 *0 ot) *1 par } }
		\ar@3[rr] ^-{k^l}
		\ar@3 [d] ^-{\hspace{2pt} k^r}
		&& { \twocell{ (1 *0 pareta *0 1) *1 (1 *0 par) *1 par} }
		\ar@[lightgray]@3 [d] ^-{\hspace{2pt} \color{lightgray}\lpc}
		\\
		{ \twocell{(1 *0 oteta *0 1) *1 (1 *0 par) *1 ot } }
		\ar@[lightgray]@3 [rr] _-{\vphantom{\rule{0pt}{1.5ex}} \color{lightgray}k^l} ^-{}="1"
		&& { \twocell{par} }
				\ar@0 "2,2"!<30pt,30pt>;"1"!<0pt,0pt>_*+{(36)}
		}\label{Crit36}
	}
}